\section{Penelitian Terdahulu}

Pendekatan evolusioner untuk ABAC \textit{policy mining} telah diperkenalkan sejak \textcite{gaspar-cunha_evolutionary_2015}.
Mereka mengusulkan optimasi multi-objektif dengan representasi dan operator genetik khusus ABAC, strategi \textit{separate-and-conquer}, inisialisasi populasi kustom, promosi diversitas, dan penghentian dini.
Tiga objektif yang dioptimalkan adalah konsistensi terhadap data akses, kompleksitas kebijakan yang rendah, dan penghindaran atribut yang secara unik mengidentifikasi subjek atau objek. 
Sebelumnya, penerapan teknik evolusioner pada kebijakan keamanan juga telah dieksplorasi secara terbatas, antara lain melalui \textit{genetic programming} untuk kebijakan multi-level \parencite{li_policy_2008} dan SPEA2\label{first:spea2} untuk pembelajaran kebijakan keamanan dinamis \parencite{lim_dynamic_2009,gaspar-cunha_evolutionary_2015}.

Perkembangan berikutnya mengarah pada metaheuristik \textit{population-based} berbasis kecerdasan kawanan. 
\textcite{shang_abac_2024} menggunakan \textit{Ant Colony Optimization} (ACO), sedangkan \textcite{li_attribute_2026} mengusulkan \textit{Upgraded Black-winged Kite} (UBK), dengan evaluasi yang mencakup akurasi dan \textit{Weighted Structural Complexity} (WSC).

Di luar pendekatan evolusioner, \textcite{guo_--fly_2025} mengusulkan kerangka \textit{on-the-fly} yang memformulasikan \textit{policy mining} sebagai optimasi aproksimasi submodular dengan pemangkasan ruang pencarian. 
Pendekatan ini berfokus pada usabilitas kebijakan bagi manusia maupun LLM\label{first:llm} dan dilaporkan mampu menurunkan redundansi sebesar 93,2\% dibandingkan metode pembanding. 
Namun, kerangka tersebut bersifat deterministik/\textit{greedy} dan ditujukan untuk kontrol akses secara umum, bukan ABAC secara spesifik.

Meskipun berbeda dalam mekanisme pencarian, \textcite{gaspar-cunha_evolutionary_2015}, \textcite{shang_abac_2024}, \textcite{li_attribute_2026}, dan \textcite{guo_--fly_2025} tidak menjadikan keterbatasan memori perangkat \textit{edge-IoT} sebagai pertimbangan utama dalam desain algoritmanya. 
Kesenjangan ini menjadi salah satu dasar pengembangan pendekatan \textit{policy mining} yang lebih hemat sumber daya pada penelitian ini.

\begin{longtable}{|c|p{2.5cm}|p{3cm}|p{3cm}|p{3cm}|}
\caption{Ringkasan Posisi Penelitian Terdahulu}
\label{tab:posisi-penelitian} \\
\hline
\textbf{No} & \textbf{Peneliti (Tahun)} & \textbf{Metode} & \textbf{Fokus/Kontribusi} & \textbf{Celah terhadap CoDA-EDA} \\
\hline
\endfirsthead

\multicolumn{5}{c}{{\bfseries \tablename\ \thetable{} -- lanjutan dari halaman sebelumnya}} \\
\hline
\textbf{No} & \textbf{Peneliti (Tahun)} & \textbf{Metode} & \textbf{Fokus/Kontribusi} & \textbf{Celah terhadap CoDA-EDA} \\
\hline
\endhead

\hline
\multicolumn{5}{r}{Lanjut ke halaman berikutnya} \\
\endfoot

\hline
\endlastfoot

% ========== DATA BARIS ==========
1 & \textcite{gaspar-cunha_evolutionary_2015} & \textit{Genetic programming} multi-objektif (representasi fenotipik kustom) & Salah satu pendekatan evolusioner paling awal untuk ABAC policy mining; 3 objektif: konsistensi, kompleksitas rendah, hindari atribut identitas unik & Population-based (menyimpan populasi pohon ekspresi penuh); tidak dirancang untuk perangkat memori terbatas \\
\hline
2 & \textcite{shang_abac_2024} & \textit{Ant Colony Optimization} (ACO) & Baseline population-based penelitian ini; dievaluasi dengan akurasi \& WSC & Population-based --- jejak memori (matriks feromon) tumbuh dengan ukuran populasi/masalah \\
\hline
3 & \textcite{li_attribute_2026} & \textit{Upgraded Black-winged Kite} (UBK) & Baseline population-based penelitian ini; kecerdasan kawanan berbasis populasi & Population-based --- kebutuhan memori sama dengan algoritma population-based lain \\
\hline
4 & \textcite{guo_--fly_2025} & Optimasi submodular/aproksimasi (deterministik, non-evolusioner) & Kontrol akses umum (bukan ABAC spesifik); fokus usabilitas kebijakan bagi manusia/LLM & Bukan pendekatan population-based/EDA; tidak membahas keterbatasan memori perangkat edge-IoT \\
\hline
5 & --- (celah penelitian) & EDA / \textit{compact-EDA} / \textit{dependency-aware} EDA & Belum ditemukan penelitian yang menerapkan kelas algoritma ini pada ABAC policy mining & Celah inilah yang diisi CoDA-EDA, bersama Premis 1-4 dan Hipotesis H1-H4 \\
\hline

\end{longtable}

Berdasarkan Tabel \ref{tab:posisi-penelitian}, penelitian terdahulu pada \textit{ABAC policy mining} belum menjadikan keterbatasan memori perangkat \textit{edge-IoT} sebagai pertimbangan desain algoritma. 
Selain itu, berdasarkan penelusuran literatur yang dilakukan, belum ditemukan penerapan \textit{Estimation of Distribution Algorithm} (EDA) pada persoalan tersebut. 
Kesenjangan ini menjadi dasar kebaruan penelitian yang diuraikan lebih lanjut pada subbab berikutnya.